%==================================================================

\subsection{Факт}
\leavevmode\par
\begin{fact}[Факт из лекции]\label{fact:hom}
Пусть $\mathcal{C}$ --- маленькая категория,
$A_1,A_2\in\mathrm{Ob}\,\mathcal{C}$,
$F_1=\Hom(A_1,-)$, $F_2=\Hom(A_2,-)$ --- функторы
$\mathcal{C}\to\mathrm{Sets}$. Если существует натуральное
преобразование $\eta\colon F_1\Rightarrow F_2$ такое, что каждая его
компонент $\eta_X\colon\Hom(A_1,X)\to\Hom(A_2,X)$ является биекцией, то
$A_1\simeq A_2$.
\end{fact}

Иными словами: функтор $\Hom(A,-)$ определяет объект категории
\emph{однозначно}. Интуиция: $A$ --- это <<тот объект, через который
удобно смотреть на все остальные>>, и если два объекта дают одинаковые
<<окна>>, они изоморфны.

\subsection{Доказательство}
\textbf{Шаг 1: достаём $f$ и $g$.}
Рассмотрим компоненту $\eta_{A_1}\colon\Hom(A_1,A_1)\to\Hom(A_2,A_1)$;
она переводит $\id_{A_1}$ в некоторый морфизм
\[
 f\colon A_2\longrightarrow A_1.
\]
Компонента $\eta_{A_2}\colon\Hom(A_1,A_2)\to\Hom(A_2,A_2)$ --- биекция,
поэтому можно рассмотреть обратную $\eta_{A_2}^{-1}$; она переводит
$\id_{A_2}$ в морфизм
\[
 g\colon A_1\longrightarrow A_2.
\]

\textbf{Шаг 2: докажем $g\circ f=\id_{A_2}$.}
Натуральность $\eta$ для морфизма $g\colon A_1\to A_2$ --- это
коммутативность диаграммы~\ref{dgm:fact1}.

\begin{diagram}{Натуральность $\eta$ для $g\colon A_1\to A_2$: стрелки
 $-\circ g$ --- композиция с $g$.\label{dgm:fact1}}
\begin{tikzcd}[cd style]
 \Hom(A_1,A_1) \ar[r,"\eta_{A_1}"] \ar[d,"g\circ-"'] &
 \Hom(A_2,A_1) \ar[d,"g\circ-"] \\
 \Hom(A_1,A_2) \ar[r,"\eta_{A_2}"'] & \Hom(A_2,A_2)
\end{tikzcd}
\end{diagram}

Пояснение к диаграмме~\ref{dgm:fact1}. Это условие натуральности
преобразования $\eta$ при морфизме $g$: маршрут <<сначала
$\eta_{A_1}$, потом скомпоновать с $g$>> (вправо-вниз) равен
маршруту <<сначала скомпоновать с $g$, потом $\eta_{A_2}$>>
(влево-вниз-вправо). Ниже оба маршрута прослеживаются на
элементе $\id_{A_1}$, и равенство даёт $g\circ f=\id_{A_2}$.

Проследим, куда попадает $\id_{A_1}$ по двум маршрутам:
\begin{bullets}
\item сверху направо: $\id_{A_1}\mapsto\eta_{A_1}(\id_{A_1})=f$, затем
 вниз: $f\mapsto g\circ f$;
\item слева вниз, затем направо: $\id_{A_1}\mapsto g$, затем
 $g\mapsto\eta_{A_2}(g)$.
\end{bullets}
Коммутативность даёт $\eta_{A_2}(g)=g\circ f$. С другой стороны, по
построению $g$ мы имели $\eta_{A_2}^{-1}(\id_{A_2})=g$, т.\,е.
$\eta_{A_2}(g)=\id_{A_2}$. Значит
\begin{equation}\label{eq:gf}
 g\circ f=\id_{A_2}.
\end{equation}

\textbf{Шаг 3: строим обратное преобразование $\eta^{-1}$.}
Положим $(\eta^{-1})_X:=\eta_X^{-1}$, т.\,е.
$\Hom(A_2,X)\to\Hom(A_1,X)$. Надо проверить натуральность: для любого
морфизма $\alpha\colon X\to Y$ диаграмма~\ref{dgm:fact2} должна быть
коммутативной.

\begin{diagram}{Натуральность обратного преобразования: компоненты
 $\eta_X^{-1}$ и композиции с $\alpha$.\label{dgm:fact2}}
\begin{tikzcd}[cd style]
 \Hom(A_2,X) \ar[r,"\alpha\circ-"] \ar[d,"\eta_X^{-1}"'] &
 \Hom(A_2,Y) \ar[d,"\eta_Y^{-1}"] \\
 \Hom(A_1,X) \ar[r,"\alpha\circ-"'] & \Hom(A_1,Y)
\end{tikzcd}
\end{diagram}

Пояснение к диаграмме~\ref{dgm:fact2}. Условие натуральности
для обратного преобразования $\eta^{-1}$. Два маршрута из
$\Hom(A_2,X)$ в $\Hom(A_1,Y)$: вправо-вниз --- «сначала
скомпоновать с $\alpha$, потом применить $\eta^{-1}$» ---
и вниз-вправо --- «сначала $\eta^{-1}$, потом композиция
с $\alpha$»; натуральность и есть их равенство. На элементе
$p$ маршруты прослеживаются ниже.

Проверка: пусть $p\colon A_2\to X$. Поскольку $\eta_X$ --- биекция,
$p=\eta_X(q)$ для единственного $q\colon A_1\to X$, т.\,е.
$q=\eta_X^{-1}(p)$. Тогда оба маршрута дают
$\eta_Y^{-1}(\alpha\circ\eta_X(q))$: слева направо --- по определению,
справа налево --- $\eta_Y^{-1}(\alpha\circ p)$. Коммутативность
доказана.

\textbf{Шаг 4: докажем $f\circ g=\id_{A_1}$.}
Рассмотрим диаграмму~\ref{dgm:fact3}, натуральность $\eta^{-1}$ для
морфизма $f\colon A_2\to A_1$, и применим её к $\id_{A_2}$.

\begin{diagram}{Натуральность $\eta^{-1}$ для $f\colon A_2\to A_1$
 (та же диаграмма \ref{dgm:fact2} с частным выбором $X=A_2$, $Y=A_1$;
 стрелки $f\circ-$).\label{dgm:fact3}}
\begin{tikzcd}[cd style]
 \Hom(A_2,A_2) \ar[r,"\eta_{A_2}^{-1}"] \ar[d,"f\circ-"'] &
 \Hom(A_1,A_2) \ar[d,"f\circ-"] \\
 \Hom(A_2,A_1) \ar[r,"\eta_{A_1}^{-1}"'] & \Hom(A_1,A_1)
\end{tikzcd}
\end{diagram}

Пояснение к диаграмме~\ref{dgm:fact3}. Частный случай
диаграммы~\ref{dgm:fact2}: naturality $\eta^{-1}$, применённая
к морфизму $f\colon A_2\to A_1$. Здесь стрелки --- композиция
с $f$; приложение коммутативности к $\id_{A_2}$ даёт
$f\circ g=\id_{A_1}$.

Верхний путь: $\id_{A_2}\mapsto\eta_{A_2}^{-1}(\id_{A_2})=g$, затем
вниз: $g\mapsto f\circ g$. Нижний путь: $\id_{A_2}\mapsto f$, затем
направо: $f\mapsto\eta_{A_1}^{-1}(f)=\eta_{A_1}^{-1}(\eta_{A_1}(\id_{A_1}))
=\id_{A_1}$. Значит
\begin{equation}\label{eq:fg}
 f\circ g=\id_{A_1}.
\end{equation}

Из~\eqref{eq:gf} и~\eqref{eq:fg} следует $A_1\simeq A_2$ (изоморфизм
$f$ с обратным $g$). \hfill$\square$

\subsection{Следствие: как узнать пучок}
\leavevmode\par
\begin{cor}\label{cor:criterion}
Пусть $\Shv{F}$ --- предпучок, $\Shv{F}'$ --- пучок. Если для всех
пучков $\Shv{G}$ изоморфизм $\Hom(\Shv{F},\Shv{G})\simeq\Hom(\Shv{F}',
\Shv{G})$ естествен (т.\,е. коммутативны все диаграммы вида
\ref{dgm:natG} и \ref{dgm:natF}), то $\Shv{F}\simeq\Shv{F}'$ как
предпучки.
\end{cor}

\begin{proof}
Возьмём $\mathcal{C}=\Sh(X)$ и $F_1=\Hom(\Shv{F},-)$,
$F_2=\Hom(\Shv{F}',-)$ как функторы $\Sh(X)\to\mathrm{Sets}$. Данное
естественное преобразование удовлетворяет условиям Факта, поэтому
$\Shv{F}\simeq\Shv{F}'$ в $\mathcal{C}$, т.\,е. как пучки; но $\Shv{F}'$
--- пучок, значит и $\Shv{F}$ изоморфен пучку.
\end{proof}

Это следствие и есть <<принцип единственности>>: два способа
<<прокачать>> предпучок, дающие одинаковые морфизмы во все пучки,
совпадают. В частности, $\Shv{F}^{\#}$ определён однозначно.

\begin{rem}[лемма Ёнеды]
Факт~\ref{fact:hom} --- частный случай знаменитого погружения
категории в категорию представимых функторов (лемма Ёнеды).
\end{rem}

\begin{bookbox}
\booksrc{Манин, \S\,1.16.8---1.16.9, с.~104---105: представимые
 функторы $h_X=\Hom(X,-)$ и теорема $f\mapsto h_f$ --- изоморфизм
 $\Hom_{\mathcal C}(X,Y)\simeq\Hom_{\mathrm{Funct}}(h_X,h_Y)$; функтор
 $h\colon\mathcal C\to\widehat{\mathcal C}$ является погружением
 категории. Именно это и есть категорная форма <<Факта>>:
 $\Hom(A,-)$ определяет объект однозначно.
 Годеман, гл.~I, п.~1.7, с.~25 --- пример того, как морфизмы
 порождают естественные преобразования функторов $\Hom$.
}
\end{bookbox}

%==================================================================
